% !PRINTED 490
\weberpage{490}
2. \gesperrt{If}
\[
P = \pi,\ \pi',\ \pi'',\ \pi''' \ldots
\]
\gesperrt{is a group of permutations of $m$ elements and $\pi_1$ is any one among them, then the system}
\[
P_1 = \pi\pi_1,\ \pi'\pi_1,\ \pi''\pi_1,\ \pi'''\pi_1 \ldots,
\]
\gesperrt{agrees completely with $P$, apart from the order.}

For the permutations contained in $P_1$, whose number is just as great as that of the permutations $P$, are, by the group property of $P$ (\S.\,148, 3.), in any case all contained among the permutations of $P$. Moreover, the permutations of $P_1$ are (by \S.\,148, 2.) all different from one another, and hence $P$ is identical with $P_1$, apart from the order. In the same way one can also show that
\[
P_1' = \pi_1\pi,\ \pi_1\pi',\ \pi_1\pi'',\ \pi_1\pi''', \ldots
\]
is identical with $P$.

From this the following theorems result:

3. \gesperrt{Every permutation group contains the identical permutation.}

For among the permutations of $P_1$, $\pi_1$ itself must also occur, and if $\pi_0\pi_1 = \pi_1$, then (by \S.\,148, 2.) $\pi_0$ is the identical permutation.

4. \gesperrt{A permutation group $P$ contains, for every permutation, also the opposite one.}

For by 3. the identical permutation must occur in $P_1$. But if $\pi\pi_1$ is the identical permutation, then $\pi = \pi_1^{-1}$. Here $\pi_1$ may be any permutation from $P$.

If in place of the variables $u, u_1 \ldots u_{m-1}$ definite quantities $\alpha, \alpha_1, \alpha_2 \ldots \alpha_{m-1}$ are put, be they numerical values or quantities of any field, then Theorem 1. loses its general validity. If we take, for example, $m = 3$ and fix between the three quantities $\alpha, \alpha_1, \alpha_2$ the relation $2\alpha_2 = \alpha + \alpha_1$
while otherwise not restricting them further, then the function $\alpha_1 - \alpha_2$ remains unchanged only under the identical permutation and under the permutation
\[
\begin{pmatrix} 0, & 1, & 2 \\ 1, & 2, & 0 \end{pmatrix}.
\]
